\section{Stage 2: sky realisation -- four modes A, B, C, E}
\label{sec:stage2}
% =====================================================================

\paragraph{Mode A. Prediction.} Sample
\(a_{\ell m}\sim\mathcal N(0,C_\ell^{\mathrm{TFPT}})\) coherently from
the TFPT block-diagonal harmonic covariance and synthesise
\(T,Q,U\) on HEALPix.  This is one TFPT universe, not our universe.
Run \(\geq 10^3\) seeds and compare \emph{distribution-level}
statistics to Planck.

\paragraph{Mode B. Conditioned reconstruction.} Compute, in harmonic
space, the Wiener filter
\(\hat a = (C^{-1}+B^\dagger N^{-1}B)^{-1} B^\dagger N^{-1} d\)
with \(C\) the TFPT spectrum, \(B\) the Planck beam/mask operator,
\(N\) the Planck noise covariance, and \(d\) the Planck map.  The
result is \emph{our} sky reconstructed under the TFPT prior;
\emph{not} a prediction.

\paragraph{Mode C. Phase-matched diagnostic.} Reuse Planck phases with
TFPT amplitudes,
\(a_{\ell m}^{\mathrm{diag}}=a_{\ell m}^{\mathrm{Planck}}
\sqrt{C_\ell^{\mathrm{TFPT}}/\widehat C_\ell^{\mathrm{Planck}}}\).
\textbf{Diagnostic only}; never reported as prediction.

\paragraph{Mode E. Lemma-E branch phase selection (NEW in v4.6).}
\label{sec:modeE_def}
The phases \(\varphi_{\ell m}\) are constructed from
\(\arg\langle e_j,u_\Sigma\rangle\) via the canonical joint
diagonalisation of the admissible primordial-curvature Hessian
\(K_\Sigma\), with the spectral flag breaking degeneracies and the
seam convention fixing the residual U(1).  No Planck phase enters; no
SMICA \(a_{\ell m}\) enters.  See~\S\ref{sec:modeE} for the full
construction and acceptance results.

\subsection{Final maps: side-by-side}
\label{sec:final_maps}

Figure~\ref{fig:mollweide} shows the four-panel comparison of Planck
SMICA, TFPT Mode A (random, lensed, \(N_{\mathrm{side}}=1024\)), TFPT
Mode E (Lemma-E, candidate v4.6, \(N_{\mathrm{side}}=512\)), and the
Planck \(-\) Mode-A residual.

\begin{figure}[!htb]
\centering
\includegraphics[width=\textwidth]{fig01_mollweide_grid.png}
\caption{TFPT v4.6.3 vs Planck PR3 SMICA. \emph{Top-left:} Planck
SMICA (PR3) downgraded to \(N_{\mathrm{side}}=1024\). \emph{Top-right:}
TFPT Mode~A (random sample from \(C_\ell^{\mathrm{TFPT,v4.6}}\) with
seed~42, lensed, birefringence~applied). \emph{Bottom-left:} TFPT
Mode~E (Lemma-E branch phases, candidate v4.6, \(N_{\mathrm{side}}=512\)).
\emph{Bottom-right:} Planck \(-\) TFPT Mode~A residual (RMS
\(\approx 153\,\mu\mathrm{K}\)).  Mode~A is statistically
indistinguishable from Planck SMICA at the eyeball level; this is
expected, because A--D give the correct \(C_\ell\) and Mode~A injects
random phases.  Mode~E uses fixed phases coming from the spectrum of
\(K_\Sigma\); the residual against Planck is \(13\,\%\) smaller.}
\label{fig:mollweide}
\end{figure}

\subsection{Residual statistics}

Figure~\ref{fig:residhist} shows the residual histograms.  Mode~A
distributes as \(\mathcal N(0,\sigma_A^2)\) with
\(\sigma_A=153.3\,\mu\mathrm{K}\); Mode~E narrows to
\(\sigma_E=133.6\,\mu\mathrm{K}\), a \(13\,\%\) reduction relative to
the same Mode-A run at the same \(N_{\mathrm{side}}\).

\begin{figure}[!htb]
\centering
\includegraphics[width=0.78\textwidth]{fig03_residual_histogram.png}
\caption{Pixel-residual histogram of Planck~\(-\)~TFPT, in
\(\mu\mathrm{K}\), restricted to the unmasked sky
(\(f_{\mathrm{sky}}\approx 0.913\)).  Mode~A: independent random
realisation, RMS~\(=153.3\,\mu\mathrm{K}\).  Mode~E: deterministic
phases from \(\arg\langle e_j,u_\Sigma\rangle\),
RMS~\(=133.6\,\mu\mathrm{K}\) (\(13\,\%\) reduction).}
\label{fig:residhist}
\end{figure}

% =====================================================================
\section{Mode E in detail: Lemma E, candidate v4.6}
\label{sec:modeE}
% =====================================================================

\subsection{The Branch Operator Bundle}

The minimal native structure of Mode~E (Stefan v4.6 brief~\S11) is
\begin{equation}
\mathfrak B_\Sigma \;=\; (\,
\mathcal H_{\mathrm{adm}},\;
K_\Sigma,\;
u_\Sigma,\;
\iota_C,\;
\mathcal O\,) ,
\end{equation}
i.e.\ the admissible Hilbert space, the primordial-curvature Hessian,
the primitive seam state, the Calderon involution, and the observer.
Implementation: \texttt{tfpt\_cmb.branch\_phase\_selection.BranchOperatorBundle}
plus \texttt{ObserverEmbedding}.

\subsection{The candidate v4.6 bundle}

In the absence of a real Sprint-2A boundary extraction, v4.6.3 ships
a deterministic \emph{candidate} bundle synthesised from the seven
primitives:
\begin{itemize}[leftmargin=1.4em]
\item \(K_\Sigma\): Hermitian, eigenvalues
\(\lambda_j = c_3^2\,j\,(1+\phi_0\cos(2\pi\beta j))\) on \(j=1\dots M\);
\item \(u_\Sigma\): seam-graded
\(u_j = \exp(2\pi i c_3 j)/\sqrt M\) modulated by
\(\cos(2\pi\lambda_C j/\Omega_{\mathrm{adm}})\), normalised;
\item \(\iota_C\): pairing involution \(j\leftrightarrow M-j+1\);
\item spectral flag: c\(_3\)-graded twist operator commuting with
\(K_\Sigma\) by construction;
\item observer: identity tetrad, \(\eta_{\mathrm{LS}}=280\,\mathrm{Mpc}\),
parity \(+1\).
\end{itemize}
This is \emph{not} a real branch realisation. It is a structural
proxy that respects all algebraic axioms and is reproducible across
runs (the rng seed only fixes the Haar-random rotation that mimics
the unknown spectral flag).  The run is tagged
\textsf{lemma\_E\_candidate\_v46\_unproven} in every JSON artifact.

\subsection{The Lemma-E phases}

Figure~\ref{fig:phij} shows the resulting branch-phase distribution
\(\{\varphi_j\}_{j=1}^{M}\) and the spectrum of \(K_\Sigma\).

\begin{figure}[!htb]
\centering
\includegraphics[width=\textwidth]{fig04_phi_j_distribution.png}
\caption{\emph{Left:} Histogram of the Lemma-E branch phases
\(\varphi_j=\arg\langle e_j,u_\Sigma\rangle\) for the \(M=192\)
candidate-v4.6 modes.  Note the explicit absence of the uniform
distribution: the seam projection imprints a non-trivial
phase structure that, in a real branch, would lift the
random-realisation degeneracy.  \emph{Right:} The spectrum
\(\{\lambda_j\}\) of the candidate \(K_\Sigma\), log-scaled.  Both
quantities are reproducible from the seven primitives, with the rng
seed fixing only the spectral-flag rotation.}
\label{fig:phij}
\end{figure}

\subsection{Acceptance checks}

Stefan v4.6 brief~\S10 lists nine acceptance checks Mode~E must pass.
Figure~\ref{fig:checks} shows the v4.6.3 status.

\begin{figure}[!htb]
\centering
\includegraphics[width=0.95\textwidth]{fig05_acceptance_checks.png}
\caption{Mode-E acceptance checks (Stefan v4.6 brief~\S10) as run
in v4.6.3 for the candidate bundle.  Seven of nine pass.  Two
remain pending: the pixel-correlation high-bar requires Lemma~E to
be \emph{proven}, not conjectured (the candidate phases happen to
align with Planck phases only structurally, not realisation-wise);
the official Planck likelihood requires a coupled foreground/noise
model that is outside Paper~7's scope.}
\label{fig:checks}
\end{figure}

\subsection{The 13\,\% RMS reduction}

The Mode-E residual RMS is \(133.6\,\mu\mathrm{K}\), i.e.\
\(13\,\%\) below the Mode-A RMS \(153.6\,\mu\mathrm{K}\) at the same
\(N_{\mathrm{side}}\).  We do \emph{not} interpret this as evidence
for Lemma~E.  In the candidate construction, the Lemma-E phases happen
to inherit a deterministic structure that is not aligned with the
Planck realisation \emph{by construction} --- the bundle is not
extracted from the boundary data.  The 13\,\% gap is consistent with
Mode~E phases being slightly less anti-correlated with Planck than a
truly random realisation would be at this \(N_{\mathrm{side}}\)
(c.\,f.~Figure~\ref{fig:residhist}).  A real Lemma-E test requires the
Sprint-2A boundary extraction.

% =====================================================================
\section{Birefringence: TFPT vs Planck EB diagnostic}
\label{sec:birefringence}
% =====================================================================

The TFPT prediction \(\beta_{\mathrm{TFPT}}=\phi_0/(4\pi)
=0.2424^\circ\) sits well inside the Eskilt et al.\ 2022 measurement
\(0.34^{+0.09}_{-0.09}{}^\circ\) and is consistent with the cosmological
birefringence reported by recent re-analyses.  The pipeline's own EB
estimator on the SMICA component map yields
\(\hat\beta_{\mathrm{EB,SMICA}}\approx 0.027^\circ\) --- a factor
\(\approx 9\) smaller than the TFPT prediction.

This is \emph{not} a contradiction.  EB on a single component-separated
map is foreground-residual and noise-dominated and does \emph{not}
recover \(\beta\); the reliable measurement requires a full
foreground-marginalised likelihood (see Eskilt+22 for the procedure).
The TFPT prediction is a derived TFPT quantity (Lemma E and
\(C_{12}\)) and is not affected by this in-pipeline EB diagnostic.
We retain the factor-9 mismatch as an open observational caveat
(\texttt{native\_closure\_v46.OBSERVATIONAL\_CAVEATS}).

% =====================================================================
\section{Comparison metrics and audit ledger}
\label{sec:metrics}
% =====================================================================

\begin{center}
\renewcommand{\arraystretch}{1.18}
\begin{tabularx}{\textwidth}{@{}p{0.27\textwidth}p{0.16\textwidth}p{0.16\textwidth}Y@{}}
\toprule
\textbf{Quantity} & \textbf{TFPT v4.6.3} & \textbf{Planck} & \textbf{Comment} \\
\midrule
\(D_{220}^{TT}\,[\mu\mathrm{K}^2]\) & 6\,546 & 6\,202 & 5.6\,\% high; tracks \(\omega_c\) excess. \\
\(C_{220}^{TT}\,[\mu\mathrm{K}^2]\)  & 0.722 & 0.758 & SMICA \texttt{anafast} \\
\(\sigma_8\) proxy (RMS T) & 44.85 & 42.21 & 6\,\% high \\
Cold-spot min. (\(\mu\mathrm{K}\))& \(-164\) & \(-173\) & 5\,\% \\
Cold-spot max. (\(\mu\mathrm{K}\))& \(+145\) & \(+150\) & 3\,\% \\
TT \(\chi^2_{\mathrm{Knox,diag}}/\)dof & 107.6 & --- & diag-only; not Likelihood \\
Mode-A residual RMS \([\mu\mathrm{K}]\) & 153.3 & --- & expected for indep.\ realisation \\
Mode-E residual RMS \([\mu\mathrm{K}]\) & 133.6 & --- & 13\,\% below Mode A \\
\(\beta_{\mathrm{birefr.}}\) (TFPT pred.)\,\([\mathrm{deg}]\) & 0.2424 & 0.34\(\pm\)0.09 (Eskilt+22) & inside \(2\sigma\) \\
\(\hat\beta\) (in-pipe EB on SMICA) & --- & 0.027 & foreground/noise dominated \\
\bottomrule
\end{tabularx}
\end{center}

\paragraph{Spectrum-shape \(\chi^2\).} A diagonal-Gaussian \(\chi^2\)
using the Knox formula
\(\sigma^2(C_\ell)=2/((2\ell+1) f_{\mathrm{sky}})\,(C_\ell+N_\ell)^2\),
binned in 25 logarithmic bins between \(\ell=2\) and \(\ell_{\max}\),
broken out for TT, TE, EE, BB.  TT diagnostic only;
TE/EE/BB/TB/EB are explicitly tagged \emph{diagnostic only} per
v4.5 rules R\(_2\) and R\(_4\), and the Knox numbers there are
foreground/noise-dominated.

\paragraph{Realisation distribution.} For each Mode-A seed we compute
six summary statistics --- \(C_\ell^{TT}|_{\ell=220}\), peak
\(D_\ell^{TT}\), \(5^\circ\)-smoothed RMS, cold-spot extreme values,
skewness, kurtosis --- and ask where the Planck value falls in the
TFPT distribution. Code: \texttt{tfpt\_cmb.compare.sky\_statistics}
and \texttt{planck\_in\_TFPT\_distribution}.

% =====================================================================
\section{Pipeline}
\label{sec:pipeline}
% =====================================================================

The pipeline is implemented in \texttt{tfpt\_cmb} with 358 passing
tests; the v4.6.3 end-to-end run is reproduced by:
\begin{lstlisting}[language=bash]
cd full-plank/theory-work

# 1. Generate the Native Closure (Lemmas A-D) -> 26-vector JSON
python3 scripts/build_native_closure_v46.py

# 2. Stage 1 + 2 (Mode A, lensed, against SMICA)
python3 scripts/run_cmb_stage.py \
  --cosmology-json output/native_closure_v46/cosmology_candidate.json \
  --planck-dir ../planck_pr3 --component smica \
  --output-dir output/cmb_stage_v463_lensed \
  --lmax 2500 --nside 1024 --seed 42 \
  --mode A --candidate-v46 --strict-map \
  --write-png --lens-potential-accuracy 4

# 3. Mode E (Lemma E, candidate v4.6 bundle)
python3 scripts/run_cmb_stage.py \
  --cosmology-json output/native_closure_v46/cosmology_candidate.json \
  --planck-dir ../planck_pr3 --component smica \
  --output-dir output/cmb_stage_v463_modeE \
  --lmax 1500 --nside 512 --seed 42 \
  --mode E --candidate-v46 --mode-E-candidate \
  --strict-map --write-png

# 4. Comparison figure suite
python3 scripts/plot_v463_comparison.py
\end{lstlisting}

The output directory \texttt{output/cmb\_stage\_v463\_summary/}
contains the figures shown in this paper plus
\texttt{summary.json} with all numeric residuals, the acceptance
status, and the reproducible command.

% =====================================================================
\section{The Sprint 2A Branch Extraction Contract}
\label{sec:sprint2a}
% =====================================================================

Stefan's v4.6.3 review converges on a single sharp statement:
\emph{the code is no longer the bottleneck; the data side and the
proof side are.}  The next sensible step is therefore neither another
plot nor another module --- it is a formal, executable
\textbf{Branch Extraction Contract} that any future Sprint-2A run must
satisfy.  This section defines that contract and reports what the
v4.6.3 candidate bundle scores against it.

\subsection{Contract specification}

\begin{tfptwarningbox}
\textbf{Sprint 2A Branch Extraction Contract}\\
\textbf{Input.} \emph{closed TFPT branch operators}.\\
\textbf{Output.}
\(K_{\zeta,\mathrm{adm}}\)\,(\texttt{npz}),
\(u_\Sigma\)\,(\texttt{npz}),
\(\iota_C\)\,(\texttt{npz}),
spectral\_flag\,(\texttt{json}+\texttt{npz}),
observer\_embedding\,(\texttt{json}),
\(k\)-grid\,(\texttt{npz}), direction grid\,(\texttt{npz}),
shell-power validation\,(\texttt{json}), branch origin\,(\texttt{json}).\\
\textbf{Hard checks.}
(1)~\(K_{\zeta,\mathrm{adm}}\) Hermitian;
(2)~\(\iota_C^2=I\);
(3)~\(\iota_C K \iota_C^\dagger=\overline{K}\);
(4)~\(\iota_C\) unitary;
(5)~\(u_\Sigma\) normalised, \([u_\Sigma]=1\);
(6)~spectral flag \emph{deterministic}, no Haar random rotation;
(7)~no Planck phase or \(a_{\ell m}\) input ever;
(8)~shell-averaged power reproduces A--D;
(9)~observer origin \(=\) ``derived'';
(10)~branch origin \(=\) ``real''.
\end{tfptwarningbox}

The contract is implemented in
\texttt{tfpt\_cmb.sprint2a\_contract} (Python) and
\texttt{scripts/validate\_sprint2a\_branch.py} (CLI).  Each hard check
maps to one of the five Lemma-E sub-theorems:
\begin{itemize}[leftmargin=1.4em]
\item \textbf{E1} (existence): required-output presence,
\(K\) Hermitian, \(\mathrm{branch\_origin}=\mathrm{real}\).
\item \textbf{E2} (canonical flag): spectral-flag derivation
``canonical'', no seed field, \([K, \mathcal F]=0\).
\item \textbf{E3} (seam uniqueness): \(u_\Sigma\) normalised,
\([u_\Sigma]=1\).
\item \textbf{E4} (reality): \(\iota_C^2=I\), unitarity,
\(\iota_C K \iota_C^\dagger=\overline{K}\).
\item \textbf{E5} (observer): tetrad in \(\mathrm{SO}(3)\),
\(\eta_{\mathrm{LS}}>0\), parity \(\in\{-1,+1\}\),
\(\mathrm{origin}=\mathrm{derived}\).
\end{itemize}

\subsection{The candidate v4.6 bundle vs the contract}

We dump the candidate v4.6 synthetic bundle and run it through the
validator:
\begin{lstlisting}[language=bash]
python3 scripts/validate_sprint2a_branch.py --candidate-v46-synthetic
\end{lstlisting}
The result is exactly the audit Stefan asked for: \(18/24\) hard
checks pass; the \emph{six} that fail are the v4.6.3 audit gap and
they are the gap by design:

\begin{center}
\renewcommand{\arraystretch}{1.13}
\begin{tabularx}{\textwidth}{@{}lp{0.16\textwidth}Y@{}}
\toprule
\textbf{Failed check} & \textbf{Sub-theorem} & \textbf{Detail} \\
\midrule
\texttt{branch\_origin\_real} & E1 & origin = `synthetic\_v46\_candidate' (must be `real') \\
\texttt{spectral\_flag\_canonical} & E2 & derivation = `haar\_random\_seeded' (must be `canonical') \\
\texttt{spectral\_flag\_no\_haar\_rotation} & E2 & seed field present in \texttt{spectral\_flag\_adm.json} \\
\texttt{iota\_C\_K\_conjugation} & E4 & \(\|\iota_C K \iota_C^\dagger - \overline{K}\|=1.7\) (must be \(<10^{-7}\)) \\
\texttt{observer\_origin\_derived} & E5 & origin = `candidate\_placeholder' (must be `derived') \\
\texttt{shell\_power\_matches\_AD} & global & \texttt{shell\_power\_validation.json} not produced \\
\bottomrule
\end{tabularx}
\end{center}

\begin{tfptfrontbox}
The 18 of 24 checks that pass --- including \(K\) Hermitian,
\(\iota_C^2=I\), \(\iota_C\) unitary, \(\|u_\Sigma\|=1\),
\([u_\Sigma]=1\), tetrad in \(\mathrm{SO}(3)\), \([K, \mathcal F]=0\),
no forbidden Planck inputs --- are the \emph{algebraic skeleton} of
Lemma~E.  The 6 that fail are the \emph{theorem content}: the canonical
flag, the \(\iota_C\)--\(K\) conjugation, the derived observer, the
real branch origin, and the shell-power consistency.  Closing those
six gaps is what would turn the v4.6.3 candidate into a v5.0 proven
Mode-E reconstruction.
\end{tfptfrontbox}

